% Part: lambda-calculus
% Chapter: lambda-definability
% Section: lists

\documentclass[../../../include/open-logic-section]{subfiles}

\begin{document}

\olfileid{lam}{ldf}{lst}
\olsection{याद्या}

यादी $[M_1, M_2, \ldots, M_n]$ पुढीलप्रमाणे परिभाषित करता येते:
\[
  [M_1, M_2, \ldots, M_n] \ident \lambd[sz][s M_1 (s M_2 (\ldots (s M_n z)))]
\]
अनौपचारिकपणे, यादीचे निरूपण करणारे पद हे त्या यादीचे
“संचायक” फलन असते. हे फलन दोन पदे स्वीकारते: पहिले
असे फलन असते, जे आतापर्यंत साचलेले मूल्य आणि नवा घटक
घेऊन नवे संचयित मूल्य देते; दुसरे म्हणजे संचयाचे आरंभीचे
मूल्य. ते एकेका घटकाचा संचय करते आणि शेवटी परिणाम देते.

\begin{digress}
  फलनीय प्रोग्रामलेखनाशी परिचय असेल, तर ही व्याख्या
  \emph{उजवी घडी} (right fold) यासारखी वाटेल:
  यादीतील घटकांवर उजवी घडी घालणारे फलन म्हणूनच
  यादी परिभाषित होते.
\end{digress}

या संकेतावर आधारलेली काही उपयुक्त फलने सहज परिभाषित
करता येतात:
\begin{align*}
  \fn{Sum} & \ident
  \lambd[l][l \, (\lambd[xa][\fn{Add} \, x \, a])\,  \num{0}]\\
  \fn{Len} & \ident
  \lambd[l][l \, (\lambd[xa][\fn{Add} \, a \, \num{1}]) \num{0}]
\end{align*}

$\fn{Sum}$ हे चर्च अंकांच्या यादीतील बेरजेची गणना
करते. ते यादीवर संचय करते: आरंभीचे मूल्य
$\num{0}$ असते आणि प्रत्येक घटकासाठी
$\fn{Add} \, x\, a$ हे नवे मूल्य देते
(येथे $x$ हा चालू घटक आणि $a$ हे संचयित मूल्य आहे).
परिणाम म्हणजे सर्व घटकांची बेरीज.

\begin{prob}
  $\fn{Len}$ काय करते? स्पष्ट करा.
\end{prob}

\begin{prob}
  यादीचा क्रम उलटवणारे फलन परिभाषित करा.
\end{prob}
\end{document}
